%% ma: L12-B09-0005
%% lop: 12
%% chuong: 3
%% bai: 9
%% dang: 12.9.3
%% loai: TN4
%% muc: TH
%% dap_an: "C"
%% nguon: "(NBV)-ĐỀ SỐ 13-MỖI NGÀY 1 ĐỀ THI 2025.pdf (D13, MỖI NGÀY 1 ĐỀ - Đề số 13), Phần I Câu 9"
%% kien_thuc: "Bài 9 (khoảng tứ phân vị của mẫu số liệu ghép nhóm)"
%% trang_thai: chinh-thuc
%% ghi_chu: "Đáp án gốc C đúng (Q1 = 4110/13, Q3 = 1355/4, hiệu khoảng 22,6)."
\begin{ex}
Kết quả khảo sát quãng đường đi được của một số chiếc xe ô tô điện của một hãng xe $V$ sau khi được sạc đầy pin được tổng kết lại ở bảng sau (đơn vị: km).
\begin{center}
\begin{tabular}{|c|c|c|c|c|c|}\hline
Quãng đường & $[280;300)$ & $[300;320)$ & $[320;340)$ & $[340;360)$ & $[360;380)$\\\hline
Số xe & $2$ & $13$ & $24$ & $8$ & $3$\\\hline
\end{tabular}
\end{center}
Khoảng tứ phân vị của mẫu số liệu trên thuộc nửa khoảng nào sau đây?
\choice{$[20;21)$}{$[21;22)$}{\True $[22;23)$}{$[23;24)$}
\loigiai{Cỡ mẫu $n=2+13+24+8+3=50$.

$\dfrac n4=12{,}5$ thuộc nhóm $[300;320)$ (số xe ở nhóm trước là $2$), nên $Q_1=300+\dfrac{12{,}5-2}{13}\cdot20=\dfrac{4110}{13}\approx316{,}15$.

$\dfrac{3n}4=37{,}5$ thuộc nhóm $[320;340)$ (số xe ở các nhóm trước là $2+13=15$), nên $Q_3=320+\dfrac{37{,}5-15}{24}\cdot20=\dfrac{1355}4=338{,}75$.

Khoảng tứ phân vị $\Delta_Q=Q_3-Q_1\approx22{,}6\in[22;23)$. Chọn C.}
\end{ex}
